% discussion_final.tex — §7 Discussion

\paragraph{Is the thinking tax obvious in hindsight?}
One might argue that the thinking tax is a tautological consequence of
output format: thinking mode must generate a reasoning chain
\emph{before} its answer, so of course it fails when given
too few tokens.
We make three responses.
First, while the \emph{mechanism} may be unsurprising \emph{post hoc},
the \emph{magnitude} is not: at budget 256, thinking mode achieves only
18.0\% accuracy on the full GSM8K set---a
\textbf{69.5\,pp} gap behind non-thinking at the same budget.
Prior to our study, the conventional wisdom held that ``more thinking is
always better''~\citep{snell2024scaling}; quantifying just how badly this
assumption can fail at finite budgets is a necessary empirical
contribution.
Second, the inverse scaling with model size
(\S\ref{sec:finding:model-size}) is genuinely counterintuitive:
larger, more capable models become \emph{worse} at thinking mode under
the same budget, not better.
Third, the natural-stop oracle
(\S\ref{sec:finding:natural-stop}) is a \emph{novel signal} that
emerges from the interaction between chain length and budget
constraints---it is not predicted by the truncation narrative alone and
has immediate practical utility for routing decisions.

\paragraph{Why does non-thinking work so well?}
For well-trained reasoning models, the knowledge required to solve grade-school math problems is sufficiently internalized that explicit chain-of-thought is unnecessary---the model can ``jump to the answer'' much as a skilled human can solve simple arithmetic without writing out each step.
The thinking mode forces the model to show its work, which at low budgets means it runs out of space before reaching the answer.

\paragraph{When \emph{should} one think?}
Our results do not argue against thinking mode \emph{in general}.
At sufficiently large budgets, thinking mode recovers and
eventually surpasses non-thinking performance.
The issue is that at practical token budgets---defined here as
$b \leq 512$ tokens, corresponding to the regime where generation
latency is below ${\sim}$2 seconds per query on consumer hardware
and output costs remain under \$0.001 per query at typical API
pricing---thinking mode's overhead
significantly outweighs its benefit.
This regime is relevant whenever latency or cost constraints dominate
(e.g., real-time chatbots, high-throughput batch processing, or
edge deployment with limited memory bandwidth).
The \method{} approach resolves this tension: think only when you must,
and default to the efficient non-thinking path otherwise.

\paragraph{Why not simply use \texttt{nothink@512}?}
On GSM8K, \texttt{nothink@512} may achieve slightly higher accuracy than
\method{} at comparable average token cost.
Two factors argue that \method{} remains the principled solution.
First, \method{} addresses a failure mode that
\texttt{nothink@512} cannot: for genuinely hard problems where
non-thinking produces an incorrect answer within budget, Stage~2's
thinking fallback provides an additional chance at the correct answer.
On the full test set ($n{=}1{,}319$), 64 of the 148 routed samples are genuine recoveries---problems where nothink was wrong but thinking produced the correct answer---yielding a net gain of 45 samples (64 recoveries $-$ 19 routing regrets; Error Analysis, \S\ref{sec:town-eval}).
Second, and most importantly, the value of \method{} increases as task difficulty rises.
On GSM8K, where non-thinking accuracy is already high (${\sim}$87.5\%), the marginal benefit of thinking fallback is modest (+3.4\,pp).
On harder benchmarks where non-thinking accuracy drops substantially (e.g., MATH-500), the recovery mechanism becomes critical, and \method{}'s advantage over pure nothink grows.
The principled contribution of \method{} is a \emph{general framework} for adaptive mode selection: it provides a floor on performance by matching each problem to the most efficient inference mode for its difficulty, with the early-stop signal as a zero-cost difficulty classifier.

\paragraph{Implications for scaling test-time compute.}
The inverse scaling of the thinking tax with model size (Finding~4, \S\ref{sec:finding:model-size}) has concerning implications for the roadmap of scaling test-time compute.
As models grow larger and more capable, their thinking chains become longer and more elaborate, exacerbating the budget-truncation problem.
This suggests that simply increasing model size will not ``solve'' the thinking tax---instead, complementary approaches (budget-aware training, adaptive routing, or non-thinking fallbacks) will become increasingly important.

\paragraph{Generalization beyond GSM8K.}
Our primary experiments focus on mathematical reasoning (GSM8K with
Qwen3-8B/27B, validated on MATH-500 with DeepSeek-R1-8B).
The thinking tax may manifest differently on tasks where
chain-of-thought is more structurally essential (e.g., multi-hop QA,
planning, code generation)---particularly tasks where non-thinking
accuracy is substantially lower.
We view systematic evaluation on harder benchmarks as an important
direction for future work.
However, the \emph{mechanism} we identify---truncation waste due to
verbose reasoning chains under constrained budgets---is fundamental
and model-agnostic.
It applies whenever (i)~the model's typical reasoning-chain length
exceeds the allocated budget, and (ii)~the answer appears at the
end of the chain rather than incrementally.
Our cross-model validation with DeepSeek-R1 confirms that the
phenomenon generalizes across both model families and benchmarks.
